\documentclass[aspectratio=169]{beamer}
\input{header}
\coursecode{JKSSB}

\title{Cybersecurity Prevention}
\subtitle{Lesson 44 --- Controls and the CIA Triad}
\author{JKSSB Computer Awareness}
\date{\today}

\begin{document}
\frame{\titlepage}

\begin{frame}{Why Prevention Comes First}
  \begin{itemize}
    \item Removing a virus after infection is costly; \key{blocking it first}
      is cheaper and safer.
    \item Prevention uses small daily habits plus a few tools that each do
      \key{one distinct job}.
    \item The exam favourite is matching: \key{which control does what}.
    \item Learn each control as \muted{name $\rightarrow$ threat it stops
      $\rightarrow$ what it cannot do}.
  \end{itemize}
  \begin{block}{The one idea}
    No single tool stops everything; layered controls protect the CIA triad.
  \end{block}
\end{frame}

\begin{frame}{The CIA Triad}
  \begin{itemize}
    \item \key{Confidentiality}: only authorised persons can read the data.
    \item \key{Integrity}: the data is accurate and has not been changed by
      unauthorised means.
    \item \key{Availability}: the data and systems are reachable when needed.
  \end{itemize}
  \begin{block}{One-line memory hook}
    Confidentiality keeps secrets secret; integrity keeps data correct;
    availability keeps systems up.
  \end{block}
\end{frame}

\begin{frame}{CIA Triad: Which Attack Breaks What}
  \begin{center}
    \small
    \begin{tabular}{p{0.20\textwidth}p{0.34\textwidth}p{0.36\textwidth}}
      \toprule
      \textbf{Goal broken} & \textbf{Example attack} & \textbf{Protective control} \\
      \midrule
      Confidentiality & Stealing passwords over open Wi-Fi & Encryption, strong passwords, MFA \\
      Integrity & Altering marks in a file & Clean backups, digital signatures, updates \\
      Availability & Ransomware locks files & Backup, antivirus, firewall \\
      \bottomrule
    \end{tabular}
  \end{center}
  \vspace{0.25cm}
  \muted{One attack can break more than one goal at once.}
\end{frame}

\begin{frame}{Antivirus: Signatures}
  \begin{itemize}
    \item \key{Antivirus} software detects, blocks, and removes malicious
      programs (malware).
    \item A \key{signature} is a stored pattern of a known virus, like a
      fingerprint on record.
    \item The antivirus compares each file against its \key{signature
      database} of known threats.
    \item New viruses appear daily, so signatures must be \key{updated
      regularly}; an outdated database misses new malware.
  \end{itemize}
  \begin{alertblock}{Trap alert}
    Signatures catch \key{known} malware. Unknown (zero-day) malware needs
    behaviour-based detection instead.
  \end{alertblock}
\end{frame}

\begin{frame}{Real-Time Scanning and Quarantine}
  \begin{itemize}
    \item \key{Real-time scanning} checks files as they are opened,
      downloaded, or copied --- before harm is done.
    \item A scheduled \key{full scan} checks the whole system on demand.
    \item \key{Quarantine} isolates a suspicious file in a safe separate
      area so it cannot run or spread.
    \item Quarantine is \key{not} deletion: the file is held safely until the
      user deletes or restores it.
  \end{itemize}
  \begin{block}{Keep the pair straight}
    Quarantine isolates; deletion removes. Quarantined files can be restored;
    deleted files cannot.
  \end{block}
\end{frame}

\begin{frame}{Firewall: The Gatekeeper}
  \begin{itemize}
    \item A \key{firewall} monitors incoming and outgoing network traffic and
      blocks what its rules forbid.
    \item It stops \key{unauthorised network access}, such as a stranger
      reaching into the office computer.
    \item A firewall can be \key{software} (on the computer) or
      \key{hardware} (on the network boundary).
    \item It does \key{not} remove viruses already on the disk --- that is the
      antivirus job.
  \end{itemize}
  \begin{alertblock}{Trap alert}
    Firewall blocks bad \key{traffic}; antivirus removes bad \key{files}. Do
    not swap them.
  \end{alertblock}
\end{frame}

\begin{frame}{Patching and Updates}
  \begin{itemize}
    \item A \key{patch} is a small fix released by the vendor for a known
      security weakness.
    \item \key{Software updates} close the holes that attackers exploit, so
      install them promptly.
    \item Update the \key{operating system, browser, and antivirus} alike;
      all three have exploitable flaws.
    \item \muted{Postponing updates leaves known doors open.}
  \end{itemize}
  \begin{exampleblock}{Office desktop example}
    The office PC shows an operating-system update notice. Installing it
    patches a known weakness; clicking ``remind me later'' for weeks leaves
    the machine exposed.
  \end{exampleblock}
\end{frame}

\begin{frame}{Strong Passwords}
  \begin{itemize}
    \item A \key{strong password} is long and mixes uppercase, lowercase,
      numbers, and symbols.
    \item Never reuse one password across sites, and never share it.
    \item Change \key{default passwords} on routers and devices at once.
    \item A \key{password manager} stores distinct passwords so the user
      remembers only one master secret.
  \end{itemize}
  \begin{block}{Weak vs strong}
    \texttt{rahul123} is weak; \texttt{RaHuL\#2026!Desk} is strong: longer,
    mixed, and unpredictable.
  \end{block}
\end{frame}

\begin{frame}{MFA: A Second Proof}
  \begin{itemize}
    \item \key{MFA} stands for Multi-Factor Authentication.
    \item It asks for \key{two or more proofs}: something you know (password)
      plus something you have (phone code) or are (fingerprint).
    \item Even if the password leaks, the attacker still lacks the
      \key{second factor}.
    \item \muted{One-time passwords (OTPs) sent to the phone are the most
      common second factor.}
  \end{itemize}
  \begin{alertblock}{Trap alert}
    A password alone is one factor. MFA is the \key{extra} check, not a
    stronger password.
  \end{alertblock}
\end{frame}

\begin{frame}{Backup: The Availability Control}
  \begin{itemize}
    \item A \key{backup} is a separate clean copy of important data, used
      when the original is lost, damaged, or locked by ransomware.
    \item It protects \key{availability}: work continues even after data loss.
    \item Keep backups \key{regular} and \key{offline or off-site}, so one
      incident cannot destroy both copies.
    \item Test restores: a backup that cannot be restored is \key{not} a
      backup.
  \end{itemize}
  \begin{block}{Keep the pair straight}
    Backup restores \key{lost} data; antivirus removes \key{malicious} data.
  \end{block}
\end{frame}

\begin{frame}{Safe Browsing}
  \begin{itemize}
    \item Look for \key{HTTPS} and the padlock before entering passwords or
      bank details.
    \item Do not click unknown links or attachments; \key{phishing} messages
      imitate trusted senders.
    \item Download software only from \key{official or trusted sources}.
    \item Log out of shared computers and avoid saving passwords on them.
  \end{itemize}
  \begin{exampleblock}{Phishing example}
    A message claims the bank blocked the account and demands an instant
    login through its link. A safe user ignores the link and opens the bank
    site by typing the address directly.
  \end{exampleblock}
\end{frame}

\begin{frame}{Secure Wi-Fi}
  \begin{itemize}
    \item Home and office Wi-Fi must use \key{WPA2 or WPA3} encryption with a
      strong password.
    \item \key{Open (password-free) public Wi-Fi} exposes traffic; attackers
      can read what passes over it.
    \item Avoid banking or official logins on open hotspots; use mobile data
      or a trusted \key{VPN} instead.
    \item Turn off auto-connect and change the router's \key{default admin
      password}.
  \end{itemize}
  \begin{alertblock}{Trap alert}
    The padlock (HTTPS) protects the website session; WPA2/WPA3 protects the
    \key{wireless link}. Both layers matter.
  \end{alertblock}
\end{frame}

\begin{frame}{Removable-Media Safety}
  \begin{itemize}
    \item Scan every \key{USB drive or memory card} with antivirus
      \key{before} opening its files.
    \item Disable \key{autorun} so a plugged-in drive cannot launch programs
      by itself.
    \item Do not plug unknown found drives into the office computer.
    \item \key{Eject} the device safely before removal to avoid corrupting
      files.
  \end{itemize}
  \begin{block}{Office desktop example}
    A visitor offers files on a pen drive. The safe steps are: scan first,
    open only the needed files, then eject --- never let it auto-run.
  \end{block}
\end{frame}

\begin{frame}{Which Control Does What}
  \begin{center}
    \begin{tabular}{ll}
      \toprule
      \textbf{Control} & \textbf{Main job} \\
      \midrule
      Antivirus + signatures & Detect and remove known malware files \\
      Firewall & Block unauthorised network traffic \\
      Patching / updates & Close known software weaknesses \\
      Strong password + MFA & Prove identity; block account theft \\
      Backup & Restore data after loss or ransomware \\
      HTTPS + safe browsing & Protect sessions; dodge phishing and bad downloads \\
      WPA2/WPA3 + VPN & Protect data on the network link \\
      Scan USB + disable autorun & Stop malware entering on drives \\
      \bottomrule
    \end{tabular}
  \end{center}
  \vspace{0.2cm}
  \muted{Read the stem's verb: remove, block, prove, restore, or protect.}
\end{frame}

\begin{frame}{Trap Alert: Controls to Keep Straight}
  \begin{itemize}
    \item Firewall blocks \key{traffic}; antivirus removes \key{files}.
    \item Quarantine \key{isolates}; deletion \key{removes}.
    \item MFA adds a \key{second proof}; a long password is still
      \key{one factor}.
    \item Backup restores \key{availability}; encryption protects
      \key{confidentiality}.
    \item HTTPS protects the \key{session}; WPA2/WPA3 protects the
      \key{Wi-Fi link}.
  \end{itemize}
\end{frame}

\begin{frame}{Lesson 44: Recall}
  \begin{itemize}
    \item CIA triad: confidentiality (secrets kept), integrity (data
      correct), availability (systems reachable).
    \item Antivirus uses signatures for known malware; real-time scanning
      checks on access; quarantine isolates.
    \item Firewall blocks unauthorised traffic; it does not clean infected
      files.
    \item Patch promptly; use long unique passwords plus MFA (Multi-Factor
      Authentication).
    \item Keep regular offline backups; browse over HTTPS, avoid open Wi-Fi
      for sensitive logins, and scan USB drives with autorun off.
  \end{itemize}
\end{frame}
\end{document}
